\clearpage\pagestyle{otMain}\fontsize{11}{15}\selectfont
\corrigesChapitre{10}{Droites du plan}
\begin{blocLivre}{cor-10-1}
\begin{corrige}{1}{Tester des points}
\demarche{Substituer les coordonnées dans l’équation cartésienne.}
Pour $A$, $2\times1-5+3=0$ : il appartient à $d$. Pour $B$, $4-6+3=1\ne0$ : il n’y appartient pas. Une équation de droite est un test exact d’appartenance ; il ne suffit pas qu’un point soit proche du tracé.
\label{fig-cor-10-1}
\figureGeo[corrige]{droite-01}
\end{corrige}
\end{blocLivre}
\begin{blocLivre}{cor-10-2}
\begin{corrige}{2}{De l’équation au tracé}
\demarche{Isoler y et choisir des coordonnées simples.}
On isole $y$ : $-2y=-3x-6$, puis $\res{y=\tfrac32x+3}$. Avec $x=0$, on obtient $(0;3)$ ; avec $y=0$, on obtient $(-2;0)$. Ces deux points distincts suffisent à tracer la droite et sont faciles à contrôler dans l’équation initiale.
\label{fig-cor-10-2}
\figureGeo[corrige]{droite-02}
\end{corrige}
\end{blocLivre}
\begin{blocLivre}{cor-10-3}
\begin{corrige}{3}{Une droite verticale}
\demarche{Reconnaître la coordonnée fixe et le déplacement parallèle possible.}
Elle contient tous les points d’abscisse 2 : c’est une verticale. Un directeur est $\res{(0;1)}$. Pour la même abscisse 2, elle contient une infinité d’ordonnées ; elle ne peut donc représenter une fonction de $x$. Une équation réduite de la forme demandée n’existe pas.
\label{fig-cor-10-3}
\figureGeo[corrige]{droite-03}
\end{corrige}
\end{blocLivre}
\begin{blocLivre}{cor-10-4}
\begin{corrige}{4}{Lire une pente}
\demarche{Calculer variation d’ordonnée divisée par variation d’abscisse.}
La pente vaut $(5-(-1))/(3-0)=6/3=\res2$. Le point $A$ est sur l’axe des ordonnées, donc l’ordonnée à l’origine vaut $-1$. L’équation est $\res{y=2x-1}$. Le point $B$ vérifie bien $5=2\times3-1$.
\label{fig-cor-10-4}
\figureGeo[corrige]{droite-04}
\end{corrige}
\end{blocLivre}
\begin{blocLivre}{cor-10-5}
\begin{corrige}{5}{Un point et une pente}
\demarche{Utiliser le point pour déterminer l’ordonnée à l’origine.}
Elle s’écrit $y=-2x+p$. En remplaçant par les coordonnées de $A$, $3=-2+p$, donc $p=5$. L’équation cherchée est $\res{y=-2x+5}$. Le point $(2;1)$ permet un second repère pour le tracé : avancer de 1 fait descendre de 2.
\label{fig-cor-10-5}
\figureGeo[corrige]{droite-05}
\end{corrige}
\end{blocLivre}
\begin{blocLivre}{cor-10-6}
\begin{corrige}{6}{Un point et un directeur}
\demarche{Annuler le déterminant entre un directeur et le déplacement depuis le point.}
Un point $M(x;y)$ est sur la droite si $\det((3;-2);(x-2;y-1))=0$. Cela donne $3(y-1)+2(x-2)=0$, soit $\res{2x+3y-7=0}$. On vérifie $A$ : $4+3-7=0$. Le point $A+(3;-2)=(5;-1)$ vérifie aussi l’équation.
\label{fig-cor-10-6}
\figureGeo[corrige]{droite-06}
\end{corrige}
\end{blocLivre}
\begin{blocLivre}{cor-10-7}
\begin{corrige}{7}{Deux points quelconques}
\demarche{Calculer la pente puis utiliser un des points pour trouver p.}
La pente est $(-2-4)/(3-(-1))=-6/4=-3/2$. Dans $4=(-3/2)(-1)+p$, on obtient $p=5/2$. Donc $\res{y=-\tfrac32x+\tfrac52}$. Une équation cartésienne équivalente est $3x+2y-5=0$. Les deux points vérifient les deux écritures.
\label{fig-cor-10-7}
\figureGeo[corrige]{droite-07}
\end{corrige}
\end{blocLivre}
\begin{blocLivre}{cor-10-8}
\begin{corrige}{8}{Parallèles ou confondues ?}
\demarche{Mettre les équations sous la même forme avant de comparer.}
La deuxième équation se réduit à $y=2x-1$ : $d_1$ et $d_2$ sont confondues. $d_3$ a la même pente 2 mais une ordonnée à l’origine différente ; elle est parallèle et distincte des deux premières. Il faut comparer les deux coefficients pour distinguer ces deux situations.
\label{fig-cor-10-8}
\figureGeo[corrige]{droite-08}
\end{corrige}
\end{blocLivre}
\begin{blocLivre}{cor-10-9}
\begin{corrige}{9}{Une intersection}
\demarche{Égaliser les ordonnées puis vérifier le couple trouvé.}
À l’intersection, les ordonnées sont égales : $2x+1=-x+7$. Donc $3x=6$, soit $x=2$. Puis $y=5$. Le point est $\res{I(2;5)}$. Contrôle dans les deux équations : $2\times2+1=5$ et $-2+7=5$. Les pentes sont distinctes, donc cette intersection est unique.
\label{fig-cor-10-9}
\figureGeo[corrige]{droite-09}
\end{corrige}
\end{blocLivre}
\begin{blocLivre}{cor-10-10}
\begin{corrige}{10}{Intersection avec une verticale}
\demarche{Utiliser directement la coordonnée imposée par la verticale.}
On remplace $x$ par $-2$ dans la seconde équation : $-6+2y=8$, donc $2y=14$ et $y=7$. L’intersection est $\res{(-2;7)}$. Aucune pente de la verticale n’est nécessaire : la substitution suffit.
\label{fig-cor-10-10}
\figureGeo[corrige]{droite-10}
\end{corrige}
\end{blocLivre}
\begin{blocLivre}{cor-10-11}
\begin{corrige}{11}{Des tarifs et des droites}
\demarche{Interpréter une intersection puis étudier le signe de la différence.}
L’égalité donne $5+2d=11+d$, soit $d=6$ ; le prix commun est 17 euros. La différence $A(d)-B(d)=d-6$ est négative avant 6 et positive après. Donc $A$ est moins cher pour $0\le d<6$, les tarifs sont égaux à 6 km, puis $B$ est moins cher pour $d>6$.
\label{fig-cor-10-11}
\figureGeo[corrige]{droite-11}
\end{corrige}
\end{blocLivre}
\begin{blocLivre}{cor-10-12}
\begin{corrige}{12}{Une famille de droites}
\demarche{Faire apparaître une abscisse qui élimine le paramètre.}
On réécrit $y=m(x-1)+2$. Pour $x=1$, le terme contenant $m$ s’annule et $y=2$. Toutes les droites passent par $\res{P(1;2)}$. Pour être parallèle à la droite donnée, il faut la pente $m=3$. On obtient $\res{d_3:y=3x-1}$, distincte de $y=3x-5$.
\label{fig-cor-10-12}
\figureGeo[corrige]{droite-12}
\end{corrige}
\end{blocLivre}
\begin{blocLivre}{cor-10-13}
\begin{corrige}{13}{Un alignement à deux méthodes}
\demarche{Changer de registre pour contrôler le résultat.}
La pente de $(AB)$ vaut $2/4=1/2$, et son équation est $y=x/2+2$. Pour $C$, $6/2+2=5$ : il appartient donc à $(AB)$. Avec les vecteurs, $\vect{AB}=(4;2)$ et $\vect{AC}=(8;4)$ ont un déterminant $4\times4-2\times8=0$. Les deux preuves établissent le même alignement.
\label{fig-cor-10-13}
\figureGeo[corrige]{droite-13}
\end{corrige}
\end{blocLivre}
\begin{blocLivre}{cor-10-14}
\begin{corrige}{14}{Un système à interpréter}
\demarche{Détecter les équations proportionnelles et vérifier les constantes.}
Dans le premier cas, la seconde équation est le double de la première : les deux droites sont confondues et le système a une infinité de solutions, tous les points de $x+2y=4$. Dans le second, la première imposerait $2x+4y=8$, incompatible avec 10 : le système n’a aucune solution. Les droites sont alors parallèles distinctes.
\label{fig-cor-10-14}
\figureGeo[corrige]{droite-14}
\end{corrige}
\end{blocLivre}
\begin{blocLivre}{cor-10-15}
\begin{corrige}{15}{Une médiane à déterminer}
\demarche{Construire le point géométrique nécessaire avant l’équation.}
Le milieu de $[BC]$ est $I(1;0)$. La droite $(AI)$ a une pente $(0-4)/(1-0)=-4$, donc $\res{y=-4x+4}$. Pour $y=2$, $2=-4x+4$, d’où $x=1/2$. L’intersection est $\res{(1/2;2)}$, qui est aussi le milieu de $[AI]$ ; cela fournit un contrôle supplémentaire.
\label{fig-cor-10-15}
\figureGeo[corrige]{droite-15}
\end{corrige}
\end{blocLivre}
